% from ega1-5-fr.tex
\oldpage[I]{134}
Le sous-préschéma de $X\times_S Y$ associé à l'immersion $\Gamma_f$
(\hyperref[I.4.2.1-fr]{4.2.1}) est appelé le \emph{graphe} du morphisme $f$~;
les sous-préschémas de $X\times_S Y$ qui sont des graphes de morphismes
$X\to Y$ se caractérisent par le fait que la restriction à un tel
sous-préschéma $G$ de la projection $p_1:X\times_S Y\to X$ est un
\emph{isomorphisme} $g$ de $G$ sur $X$~: $G$ est alors le graphe du morphisme
$p_2\circ g^{-1}$, où $p_2$ est la projection $X\times_S Y\to Y$.

Lorsqu'on prend en particulier $X=S$, les $S$-morphismes $S\to Y$, qui ne
sont autres que les $S$-sections de $Y$ (\hyperref[I.2.5.5-fr]{2.5.5}) sont
égaux à leurs morphismes graphes~; les sous-préschémas de $Y$ qui sont des
graphes de $S$-sections (autrement dit, ceux qui sont isomorphes à $S$ par la
restriction du morphisme structural $Y\to S$) s'appellent encore les
\emph{images} de ces sections, ou, par abus de langage, les
\emph{$S$-sections} de $Y$.

\begin{corollary}[5.3.12]
\label{I.5.3.12-fr}
Les hypothèses et notations étant celles de
(\hyperref[I.5.3.11-fr]{5.3.11}), pour tout morphisme $g:S'\to S$, soit $f'$
l'image réciproque de $f$ par $g$ (\hyperref[I.3.3.7-fr]{3.3.7})~; alors
$\Gamma_{f'}$ est l'image réciproque de $\Gamma_f$ par $g$.
\end{corollary}

C'est un cas particulier de la formule
(\hyperref[I.3.3.10.1-fr]{3.3.10.1}).

\begin{corollary}[5.3.13]
\label{I.5.3.13-fr}
Soient $f:X\to Y$, $g:Y\to Z$ deux morphismes~; si $g\circ f$ est une
immersion (resp. une immersion locale), il en est de même de $f$.
\end{corollary}

En effet, $f$ se factorise en
\[
  X\xrightarrow{\Gamma_f}X\times_Z Y\xrightarrow{p_2}Y.
\]
D'autre part, $p_2$ s'identifie à $(g\circ f)\times_Z 1_Y$
(\hyperref[I.3.3.4-fr]{3.3.4})~; si $g\circ f$ est une immersion (resp. une
immersion locale), il en est de même de $p_2$
(\hyperref[I.4.3.1-fr]{4.3.1} et \hyperref[I.4.5.5-fr]{4.5.5}), et comme
$\Gamma_f$ est une immersion (\hyperref[I.5.3.11-fr]{5.3.11}), on conclut
par (\hyperref[I.4.2.4-fr]{4.2.4}) (resp.
(\hyperref[I.4.5.5-fr]{4.5.5})).

\begin{corollary}[5.3.14]
\label{I.5.3.14-fr}
Soient $j:X\to Y$, $g:X\to Z$ deux $S$-morphismes. Si $j$ est une immersion
(resp. une immersion locale), il en est de même de $(j,g)_S$.
\end{corollary}

En effet, si $p:Y\times_S Z\to Y$ est la première projection, on a
$j=p\circ(j,g)_S$, et il suffit d'appliquer
(\hyperref[I.5.3.13-fr]{5.3.13}).

\begin{proposition}[5.3.15]
\label{I.5.3.15-fr}
Si $f:X\to Y$ est un $S$-morphisme, le diagramme
\[
  \xymatrix{
    X\ar[r]^{\Delta_X}\ar[d]_f &
    X\times_S X\ar[d]^{f\times_S f}\\
    Y\ar[r]_{\Delta_Y} &
    Y\times_S Y
  }
  \tag{5.3.15.1}\label{I.5.3.15.1-fr}
\]
est commutatif (en d'autres termes $\Delta_X$ est un morphisme fonctoriel
dans la catégorie des préschémas).
\end{proposition}

La vérification est immédiate et laissée au lecteur.

\begin{corollary}[5.3.16]
\label{I.5.3.16-fr}
Si $X$ est un sous-préschéma de $Y$, la diagonale $\Delta_X(X)$ s'identifie à
un sous-préschéma de $\Delta_Y(Y)$, dont l'espace sous-jacent s'identifie à
\[
  \Delta_Y(Y)\cap p_1^{-1}(X)=\Delta_Y(Y)\cap p_2^{-1}(X)
\]
($p_1$, $p_2$ projections de $Y\times_S Y$).
\end{corollary}

Appliquons (\hyperref[I.5.3.15-fr]{5.3.15}) au morphisme d'injection
$f:X\to Y$~; on sait alors que $f\times_S f$ est une immersion, identifiant
l'espace sous-jacent à $X\times_S X$ au sous-espace
$p_1^{-1}(X)\cap p_2^{-1}(X)$ de $Y\times_S Y$
(\hyperref[I.4.3.1-fr]{4.3.1})~; en outre, si
$z\in\Delta_Y(Y)\cap p_1^{-1}(X)$, on a $z=\Delta_Y(y)$

